%*******************************************************
% Abstract
%*******************************************************
\pdfbookmark[0]{Zusammenfassung}{Zusammenfassung}
\chapter*{Zusammenfassung}

Diese Arbeit beschäftigt sich mit der Ausführung von modellierten Geschäftsprozessen auf IoT-Systemen, um die Lücke zwischen Modellierung und Ausführung zu schließen, sodass eine prozessgetreuere Produktion bewerkstelligt werden kann. Dabei werden Prozesse im Industriebereich betrachtet, die anhand der Prozessnotation einer Ereignisgesteuerten Prozesskette (EPK) modelliert wurden. Diese Modellierungen werden mit einem eigenen Compiler in eine ausführbare Form überführt, welche auf Internet-of-Things (IoT) Systemen ausgeführt werden kann. Dabei werden formale Grundlagen gelegt bezüglich
der Geschäftsprozessmodellierung, der theoretischen Informatik sowie eingebetteten Systemen und dem Internet of Things.
Nach der Definition dieser Grundlagen findet die Umsetzung durch die Modifizierung der Prozessplattform ARIS statt, welche anschließend verwendet wird,
um einen modellierten Produktionsprozess auf IoT Systemen auszuführen. Auf Basis der Durchführung, der Erkenntnisse aus der Entwicklung und Literatur zur Thematik wird am Ende das Fazit gezogen, inwiefern solch ein Vorgehen sinnvoll ist und welche Vor- und Nachteile die Orchestrierung und Ausführung 
von Geschäftsprozessen auf IoT Systemen mit sich bringt.
